\documentclass[10pt,a4paper]{amsart}
\usepackage[margin=19mm]{geometry}
\usepackage{amsmath,amssymb,mathtools}
\usepackage[scr=rsfs,cal=euler]{mathalfa}
\usepackage{xfrac}
\usepackage[unicode,hidelinks,bookmarks=false]{hyperref}
\newcommand{\ie}{i.e.\ }
\newcommand{\cf}{cf.\ }
\newcommand{\resp}{respectively\ }
\newcommand{\Subsection}{\subsection}
\providecommand{\nobreakdash}{\nobreak\hbox{-}}
\setlength{\emergencystretch}{2em}
\multlinegap0pt
\usepackage{bm}
\usepackage{bbm}
\usepackage{dsfont}
\usepackage{xspace}
\usepackage{cancel}
\usepackage{mathtools}
\usepackage{amsthm}
\makeatletter
\@ifundefined{theorem}{\newtheorem{theorem}{Theorem}}{}
\makeatother
\title[Does Privacy Always Harm Fairness? Data-Dependent Trade-offs]{Does Privacy Always Harm Fairness? Data-Dependent Trade-offs via Chernoff Information Neural Estimation}
\author{Arjun Nichani, Hsiang Hsu, Chun-Fu (Richard) Chen, Haewon Jeong}
\date{Source: arXiv:2601.13698v2 (CC BY 4.0)}
\begin{document}
\maketitle

\noindent\emph{Source and licence.} Does Privacy Always Harm Fairness? Data-Dependent Trade-offs via Chernoff Information Neural Estimation. Version arXiv:2601.13698v2 (CC BY 4.0), distributed under the Creative Commons Attribution 4.0 International licence, as recorded in the arXiv licence metadata for this exact version and shown on the abstract page https://arxiv.org/abs/2601.13698v2. Full rights notice: licenses/arxiv-2601.13698-CC-BY-4.0.txt.\par\smallskip

\noindent\textit{Editorial scope.} Counted unit: the Theorem whose source label is \texttt{None} in the article (source0.tex lines 1107--1120, source label \texttt{None}), delivered together with the article's own complete original proof of it (source0.tex lines 1122--1214, one \texttt{proof} environment of 93 lines). No mathematics is written, repaired or re-derived: statement and proof are the article's own text line for line. Required context retained uncounted: thm:central (lines 328--336). Every definition, display and label of those lines is retained unchanged. Omitted from this package and not redistributed: 1--327, 337--1106, 1121--1121, 1215--1772. Nothing inside a retained excerpt is omitted or altered: every retained body is delivered exactly as the article's lines are written, so no omission receipt of any kind accompanies this package. Citations: the retained excerpts carry no citation command at all, so no bibliography entry of the article is transcribed and no reference is redistributed. Reference adaptation: none was needed and none was made; every reference inside the delivered unit resolves inside it, so no unresolved reference or citation is produced. Printed slips kept literal: no printed word, symbol, hypothesis or number of the counted statement or of its proof was altered. Layout and class: the article's own class is \texttt{article} with packages including subcaption, geometry, natbib, xcolor, graphicx, amsmath, amssymb, amsthm, enumitem, algorithm, algpseudocode, bm, comment, authblk; the wrapper is the lane's pinned \texttt{amsart} editorial wrapper at 10pt on A4, which restates the theorem-like environments the retained text needs and adds the article's own macro definitions that the retained text needs (none), with extra wrapper packages (bm, bbm, dsfont, xspace, cancel, mathtools). The delivered numbering is the unit's own. Assets: no retained span contains a figure environment, a graphics-inclusion command, a PDF-page inclusion, a table or a TikZ picture; no image file is copied into this package, the package declares no asset dependency and no image file is redistributed with it. Licence: version arXiv:2601.13698v2 of this article is distributed under the Creative Commons Attribution 4.0 International licence, as recorded in the arXiv licence metadata for this exact version and shown on the abstract page https://arxiv.org/abs/2601.13698v2. A full rights notice is delivered at licenses/arxiv-2601.13698-CC-BY-4.0.txt. Characters: every retained excerpt is pure ASCII, so the delivered engine, XeLaTeX with its default Latin Modern text font, reports no missing character.
\begin{theorem} \label{thm:central}
Suppose $P_0(x)\sim\mathcal{N}(\mu_0,\sigma^2\mathbf{I})$, $P_1(x)\sim\mathcal{N}(\mu_1,\sigma^2\mathbf{I})$, $Q_0(x)\sim\mathcal{N}(\zeta_0,\tau^2\mathbf{I})$, and $Q_1(x)\sim\mathcal{N}(\zeta_1,\tau^2\mathbf{I})$. Without loss of generality, we assume that $\|\mu_0 - \mu_1\|_2 \geq \|\zeta_0 - \zeta_1\|_2$. There are three behaviors of the Noisy Chernoff Difference ($\widetilde{\text{CD}}_{\eta^2}$): (i) $\widetilde{\text{CD}}_{\eta^2}$ has a maximum point, (ii) $\widetilde{\text{CD}}_{\eta^2}$ has a maximum point \underline{and} a reflection point (where $\widetilde{\text{CD}}_{\eta^2} = 0$), (iii) $\widetilde{\text{CD}}_{\eta^2}$ is non-increasing.\footnote{When $\|\mu_0 - \mu_1\|_2 = \|\zeta_0 - \zeta_1\|_2$, $\widetilde{\text{CD}}_{\eta^2}$ will always fall into this case.} The conditions for these three cases are given as follows:
    %\vspace{-1mm}
    \begin{enumerate}
        \item[(i)] $\frac{\|\zeta_0 - \zeta_1\|^2_2}{\|\mu_0 - \mu_1\|^2_2} <\frac{\tau^2}{\sigma^2} < \frac{\|\zeta_0 - \zeta_1\|_2}{\|\mu_0 - \mu_1\|_2} < 1$,  \hfill \textbf{(Case 1)}
        \item [(ii)] $\frac{\tau^2}{\sigma^2} < \frac{\|\zeta_0 - \zeta_1\|^2_2}{\|\mu_0 - \mu_1\|^2_2} < 1$, \hfill \textbf{(Case 2)}
        \item [(iii)] Neither condition (i) or (ii) hold. \hfill \textbf{(Case 3)}
    \end{enumerate}    %\vspace{-2mm}
\end{theorem}

\begin{theorem}[ \textbf{Restatement of Theorem \ref{thm:central}}]
    Suppose $P_0(x)\sim\mathcal{N}(\mu_0,\sigma^2\mathbf{I})$, $P_1(x)\sim\mathcal{N}(\mu_1,\sigma^2\mathbf{I})$, $Q_0(x)\sim\mathcal{N}(\zeta_0,\tau^2\mathbf{I})$, and $Q_1(x)\sim\mathcal{N}(\zeta_1,\tau^2\mathbf{I})$. Without loss of generality, we assume that $\|\mu_0 - \mu_1\|_2 \geq \|\zeta_0 - \zeta_1\|_2$. There are three behaviors of the Noisy Chernoff Difference ($\widetilde{\text{CD}}_{\eta^2}$): (i) $\widetilde{\text{CD}}_{\eta^2}$ has a maximum point, (ii) $\widetilde{\text{CD}}_{\eta^2}$ has a maximum point \underline{and} a reflection point (where $\widetilde{\text{CD}}_{\eta^2} = 0$), (iii) $\widetilde{\text{CD}}_{\eta^2}$ is non-increasing.\footnote{When $\|\mu_0 - \mu_1\|_2 = \|\zeta_0 - \zeta_1\|_2$, $\widetilde{\text{CD}}_{\eta^2}$ will always fall into this case.} The respective conditions for these three cases are given as follows:
    %\vspace{-1mm}
    \begin{enumerate}
    \item[]
        \item[(i)] $\frac{\|\zeta_0 - \zeta_1\|^2_2}{\|\mu_0 - \mu_1\|^2_2} <\frac{\tau^2}{\sigma^2} < \frac{\|\zeta_0 - \zeta_1\|_2}{\|\mu_0 - \mu_1\|_2} < 1$,  \hfill \textbf{(Case 1: Maximum Point)}

    \item[]     
        \item [(ii)] $\frac{\tau^2}{\sigma^2} < \frac{\|\zeta_0 - \zeta_1\|^2_2}{\|\mu_0 - \mu_1\|^2_2} < 1$, \hfill \textbf{(Case 2: Maximum and Reflection)}
        
    \item[]  
        \item [(iii)] Neither condition (i) or (ii) hold. \hfill \textbf{(Case 3: Non-increasing)}
    \end{enumerate}    %\vspace{-2mm} 
\end{theorem}

\begin{proof}

 
 Recall the definition of noisy Chernoff Difference. We define a signed noisy Chernoff Difference $s\widetilde{\text{CD}}_{\eta^2}$ such that $|s\widetilde{\text{CD}}_{\eta^2}| = \widetilde{\text{CD}}_{\eta^2}$. Thus,

\begin{align}
      s \widetilde{\text{CD}}_{\eta^2}=
     \frac{1}{8(\tau^2 + \eta^2)}\|\zeta_0 - \zeta_1\|^2_2 - \frac{1}{8(\sigma^2+\eta^2)}\|\mu_0 - \mu_1\|^2_2.
\end{align}

\paragraph{Case 1 and 2.}Let $p = \|\mu_0 - \mu_1\|_2$ and let $q = \|\zeta_0 - \zeta_1\|_2$. First, we can analyze the positive $\eta^2$ regime for a critical point. To find potential critical points, consider the derivative of $s\widetilde{\text{CD}}_{\eta^2}$.

\begin{align}
    \frac{\partial}{\partial\eta^2} s \widetilde{\text{CD}}_{\eta^2}  = \frac{p^2}{8(\sigma^2 + \eta^2)^2} - \frac{q^2}{8(\tau^2 + \eta^2)^2} = \frac{p^2(\tau^2 + \eta^2)^2-q^2(\sigma^2 + \eta^2)^2}{8(\sigma^2 + \eta^2)^2(\tau^2 + \eta^2)^2}
\end{align}

Now, the potential critical points will be $\eta^2$ where $s\widetilde{\text{CD}}_{\eta^2}' = 0$. That is $0 = p^2(\tau^2 + \eta^2)^2-q^2(\sigma^2 + \eta^2)^2$. Thus, the critical points are:

\begin{align}
    \eta_0^2 = \frac{\sigma^2q-\tau^2p}{p-q} \\\eta_1^2 =\frac{-\sigma^2q-\tau^2p}{p+q}
\end{align}



However, we observe that $\eta_1^2$ is always negative, thus the only useful critical point in the postive $\eta^2$ region is $\eta_0^2$. By our assumption, we know that $p-q \geq 0$. First, we observe that there is no critical point when $p=q$ as $\eta_0^2$ does not exist. Thus, $\eta_0^2$ is positive when the following conditions hold.

\begin{align}
    \sigma^2q-\tau^2p >0 \\
    \frac{\tau^2}{\sigma^2} < \frac{\|\zeta_0 - \zeta_1\|_2}{\|\mu_0 - \mu_1\|_2} < 1
\end{align}

Next, we will show that this critical point is always a maximum in the positive $\eta^2$ regime.

First, consider the second derivative of the signed noisy Chernoff difference.

\begin{align}
    \frac{\partial^2}{\partial (\eta^2)^2}s\widetilde{\text{CD}}_{\eta^2} =  \frac{q^2}{4(\tau^2 + \eta^2)^3} - \frac{p^2}{4(\sigma^2 + \eta^2)^3}.
\end{align}

Plugging in the relevant critical point we observe that

\begin{align}
    \frac{\partial^2}{\partial (\eta^2)^2 }s  \widetilde{\text{CD}}_{\eta_0^2} &=  \frac{q^2}{4(\tau^2 + \eta_0^2)^3} - \frac{p^2}{4(\sigma^2 + \eta_0^2)^3} = \frac{q^2}{4(\frac{q(\sigma^2-\tau^2)}{p-q})^3} -  \frac{p^2}{4(\frac{p(\sigma^2-\tau^2)}{p-q})^3} \\&=  \frac{q^2(p-q)^3}{4q^3(\sigma^2-\tau^2)^3} - \frac{p^2(p-q)^3}{4p^3(\sigma^2-\tau^2)^3} = \frac{(p-q)^4}{4p^3q^3(\sigma^2-\tau^2)^3}.
\end{align}
Now, we know that $\sigma^2 > \tau^2$ so the second derivative of  this critical point of $s\widetilde{\text{CD}}_{\eta^2}$ must be a minimum. However, the goal is to examine the behavior of $\widetilde{\text{CD}}_{\eta^2}$. So, we can show that this is critical point is a maximum of $\widetilde{\text{CD}}_{\eta^2}$, by showing that $s\widetilde{\text{CD}}_{\eta_0^2}$ is negative. By plugging in the critical point, we can observe that it is a maximum of $\widetilde{\text{CD}}_{\eta^2}$.

\begin{align}
    s\widetilde{\text{CD}}_{\eta_0^2} &= \frac{q^2}{8(\tau^2 + \eta_0^2)} - \frac{p^2}{8(\sigma^2 + \eta_0^2)} = \frac{q^2}{8(\tau^2 + \frac{\sigma^2q-\tau^2p}{p-q})} - \frac{p^2}{8(\sigma^2 + \frac{\sigma^2q-\tau^2p}{p-q})} \\&= \frac{(p-q)}{8}(\frac{q}{\sigma^2-\tau^2}-\frac{p}{\sigma^2-\tau^2}) = \frac{(p-q)(q-p)}{8(\sigma^2 - \tau^2)}
\end{align}


Now, this value is always negative as we know $p > q$ and $\sigma^2 > \tau^2$. Thus, the positive critical point is a maximum. Now, we can analyze the positive $\eta^2$ regime for a reflection point. 

\begin{align}
    &s\widetilde{\text{CD}}_{\eta^2} = \frac{q^2}{8(\tau^2 + \eta^2)} - \frac{p^2}{8(\sigma^2 + \eta^2)} \\ \notag
    \\
    &8q^2(\sigma^2 + \eta^2) - 8p^2(\tau^2 + \eta^2) = 0 \\ \notag
    \\
    &\eta^2 = \frac{q^2\sigma^2-p^2\tau^2}{p^2-q^2} = \frac{\|\zeta_0 - \zeta_1\|_2^2\sigma^2-\|\mu_0 - \mu_1\|_2^2\tau^2}{\|\mu_0 - \mu_1\|_2^2-\|\zeta_0 - \zeta_1\|_2^2}.
\end{align}





Now, the denominator of this is always positive, thus, $\eta^2$ is positive when the following holds:

\begin{align}
    \|\zeta_0 - \zeta_1\|_2^2\sigma^2-|\mu_0 - \mu_1\|_2^2\tau^2 > 0
    \\ \frac{\tau^2}{\sigma^2} < \frac{\|\zeta_0 - \zeta_1\|^2_2}{\|\mu_0 - \mu_1\|^2_2} < 1.
\end{align}

\paragraph{Case 3.} Finally, we can examine the behavior when none of these conditions hold. Suppose $\frac{\tau^2}{\sigma^2} \geq \frac{\|\zeta_0 - \zeta_1\|_2}{\|\mu_0 - \mu_1\|_2} = \frac{q}{p}$. We can analyze the sign of the numerator of $s\widetilde{\text{CD}}_{\eta^2}'$ by writing it as

\begin{align}
    &p^2(\tau^2 + \eta^2)^2 - q^2(\sigma^2 + \eta^2)^2 \\
            &= (p^2(\tau^2 + \eta^2) - q^2(\sigma^2 + \eta^2))(p^2(\tau^2 + \eta^2) + q^2(\sigma^2 + \eta^2)).
\end{align}

Thus, to analyze the sign, we can analyze $p^2(\tau^2 + \eta^2) - q^2(\sigma^2 + \eta^2)$. We observe

\begin{align}
    p^2(\tau^2 + \eta^2) - q^2(\sigma^2 + \eta^2) = p^2\tau^2 - q^2\sigma^2 + \eta^2(p^2 - q^2).
\end{align}

Now, from our initial assumption, $p^2 - q^2 \geq 0$. From the assumption that $\frac{\tau^2}{\sigma^2} \geq  \frac{q}{p}$, $p^2\tau^2 - q^2\sigma^2 \geq 0$. Thus, the sign is always positive. Now, to show that $\widetilde{\text{CD}}_{\eta^2}$ is non-increasing, we will show $s\widetilde{\text{CD}}_{\eta^2}$ is not positive. 

\begin{align}
    s\widetilde{\text{CD}}_{\eta^2} &= \frac{q^2}{8(\tau^2 + \eta^2)} - \frac{p^2}{8(\sigma^2 + \eta^2)}  = \frac{q^2(\sigma^2 + \eta^2) - p^2(\tau^2 + \eta^2)}{8((\tau^2 + \eta^2))(\sigma^2 + \eta^2)}
\end{align}

Now, we analyze $q^2(\sigma^2 + \eta^2) - p^2(\tau^2 + \eta^2)$. We can see that this is equivalent to $q^2\sigma^2 - p^2\tau^2 + \eta^2(q^2 - p^2)$. From our initial assumption, $q^2 - p^2 \leq 0$. From the assumption that $\frac{\tau^2}{\sigma^2} \geq \frac{q}{p}$, $q^2\sigma^2 - p^2\tau^2 \leq 0$. Thus, $s\widetilde{\text{CD}}_{\eta^2}$ is always negative and $\widetilde{\text{CD}}_{\eta^2}$ is non-increasing over the positive $\eta^2$ regime.
\end{proof}

\end{document}
